% =====================================================================
%  第 5 章（续）：5.4–5.6
%  原文 PDF p.154–p.176（印刷 p.141–163）
%  说明：PDF p.154 上半页仍是 5.3.3 节的结尾（该节起于 p.152，属 U07 页段），
%        为不遗漏译文，该页 5.3.3 节剩余的 5 个自然段保留在本文件开头；本文件不含 \chapter。
% =====================================================================
% [合并接缝 · 已由 Lead 处理完毕] 已删除 ch05-1.tex 末尾的 \end{proof} 与本文件开头的
% \begin{proof}[证明（续）]；5.3.3 节证明在 ch05-1 中开启、在本段末 \end{proof} 关闭，
% 合并后为单一连续证明环境。请勿再改动此接缝。

分别由 $\cA^{\fg}$ 与 $\operatorname{Vect} L\cU$ 中元素的有限线性组合表示，由此可知这个商映射是李代数同态。至于“我们的映射也是李代数同态”这一断言，则可由 $\tilde{\cA}^{\fg}_{\le1}$ 与 $\operatorname{Vect} L\cU$ 上的李括号关于各自理想所定义拓扑的连续性推出。
\end{proof}

因此，由命题 5.3.1 得到：$\tilde{\cA}^{\fg}_{\le1}$ 是李代数 $\operatorname{Vect} L\cU$ 通过其模 $\tilde{\cA}^{\fg}_0 = \operatorname{Fun} L\cU$ 的一个扩张。

然而，这个扩张与定义通常的 $L\cU$ 上阶数不超过 1 的微分算子李代数的标准（分裂）扩张并不相同（它对应 [Frenkel and Ben-Zvi 2004] 第 12.1.3 节中定义的微分算子的另一种完备化）。原因在于，代数 $\tilde{\cA}^{\fg}$ 并不作用在 $L\cU = \cU((t))$ 上的函数空间上（详见 [Frenkel and Ben-Zvi 2004] 第 12 章），而是作用在下面定义的模 $M_{\fg}$ 上。这个模可以看成支撑在子空间 $L_+\cU = \cU[[t]]$ 上的“delta 函数”空间。也正因如此，有限维情形中用过的技巧——把一个向量场提升为湮灭常数函数 1 的微分算子——不再奏效：模 $M_{\fg}$ 中没有被所有 $a_{\alpha,n}$ 湮灭的函数！所以短正合列 (5.3-7) 没有显然的分裂。事实上，下文将看到它不可分裂。\transnote{卡点：为什么这里不能像有限维情形那样，把向量场提升为湮灭常数函数 $1$ 的微分算子。背景：有限维时 $G/B$ 上的常函数被所有向量场湮灭，于是 $\cA^{\fg}_0$ 含一个一维不变子空间，提升才可行；圈情形中 $\tilde\cA^{\fg}$ 只作用在 $M_{\fg}$ 上，而 $M_{\fg}$ 是由 $|0\rangle$ 生成、被 $L_+\fn_+$ 湮灭的“delta 函数”模，其中除 $0$ 外没有向量被全部 $a_{\alpha,n}$（$n\in\Z$）湮灭。结论：扩张 (5.3-7) 的非分裂性来自模 $M_{\fg}$ 的几何性质，而非技术性障碍。[来源：原书 p.141；Frenkel--Ben-Zvi, \emph{Vertex Algebras and Algebraic Curves}, 2nd ed., AMS 2004, Ch.~12（原书此处所引）。]}

% [Tutor Note: 这里是全书“临界水平非平凡性”的几何根源。关键对比：有限维时 g 作用在 G/B 上的向量场可提升为微分算子，因为常函数提供了一维不变子空间；圈情形中 Ã^g 只作用在“delta 函数模” M_g 上，M_g 没有真空之外的不动向量，故 (5.3-7) 天然非分裂。第 6 章会看到，正是这个上同调障碍迫使 Wakimoto 模的构造必须引入额外的自由场。]

因此，不能指望像有限维情形那样把同态 $\widehat{\rho}: L\fg \to \operatorname{Vect} L\cU$ 提升为同态 $L\fg \to \tilde{\cA}^{\fg}_{\le1}$。不过我们将证明，$\widehat{\rho}$ 可以提升为同态
\[
\gghat_{\kappa_c} \to \tilde{\cA}^{\fg}_{\le1}, \qquad \mathbf{1}\mapsto 1,
\]
其中 $\gghat_{\kappa_c}$ 是 1.3.6 节定义的临界水平仿射 Kac--Moody 代数。

我们的近期目标就是证明这一断言。下一节先证明：嵌入 $L\fg \to \operatorname{Vect} L\cU$ 的像落在“局部”向量场构成的李子代数 $\cT^{\fg}_{\mathrm{loc}} \subset \operatorname{Vect} L\cU$ 之中。这一观察使我们能把扩张 (5.3-7) 替换成它的“局部”部分，后者小得多，因而更便于处理。

\section{顶点代数解释}

\begin{keypoints}
    \item \textbf{核心思想：}把 (5.3-7) 中的巨大拓扑代数换成由顶点算子的 Fourier 系数生成的“局部”部分，从而把扩张问题压缩为可计算的对象。
    \item \textbf{数学工具：}Heisenberg 顶点代数 $M_{\fg}$、自由场 $a_\alpha(z)$ 与 $a^*_\alpha(z)$、正规序、Wick 公式。
    \item \textbf{关键结论：}$\fg((t))$ 的像落在局部向量场 $\cT^{\fg}_{\mathrm{loc}}$ 中，于是 (5.3-7) 限制为局部扩张 (5.4-3)，它仍不可分裂。
    \item \textbf{记号提醒：}$\cA^{\fg}$ 是 $\cU((t))$ 上微分算子代数的一种完备化；$M_{\fg}$ 不是函数空间，而是支撑在 $L_+\cU$ 上的 delta 函数模。
\end{keypoints}

把注意力限制在短正合列 (5.3-7) 中出现的李代数的更小“局部”子代数上，会带来方便。这些更小的子代数对我们的目的已经够用，因为下文将看到，$\fg((t))$ 的像落在 $\operatorname{Vect} L\cU$ 的“局部”版本之中。为了定义这些子代数，我们把一切改用顶点代数的语言重述。

\subsection{Heisenberg 顶点代数}

设 $M_{\fg}$ 是 $\cA^{\fg}$ 的 Fock 表示，由向量 $|0\rangle$ 生成，并满足
\[
a_{\alpha,n}|0\rangle = 0,\quad n\ge 0;\qquad a^*_{\alpha,n}|0\rangle = 0,\quad n>0.
\]
显然，$\cA^{\fg}$ 在 $M_{\fg}$ 上的作用可延拓为拓扑代数 $\tilde{\cA}^{\fg}$ 的连续作用（这里给 $M_{\fg}$ 赋离散拓扑）。此外，$M_{\fg}$ 带有如下定义的 $\mathbb{Z}_+$-分次顶点代数结构（顶点代数的定义见 2.2.2 节；我们借助重构定理 2.2.4 来验证下面引入的结构满足顶点代数公理）：

\begin{itemize}
    \item $\mathbb{Z}_+$-分次：$\deg a_{\alpha,n} = \deg a^*_{\alpha,n} = -n$，$\deg |0\rangle = 0$；\transnote{卡点：$\deg=-n$ 与后两条对易关系里的负号为何必须成套出现。背景：顶点代数的平移算子 $T$ 满足 $Y(TA,z)=\partial_z Y(A,z)$，它把次数提高 $1$，即 $\deg T=1$；由 $\deg a_{\alpha,n}=-n$ 直接得到 $[T,a_{\alpha,n}]=-n\,a_{\alpha,n-1}$。而 $a^*$ 的场取 $\sum_n a^*_{\alpha,n}z^{-n}$（比 $a$ 的 $\sum_n a_{\alpha,n}z^{-n-1}$ 少一位），故其系数位移一位，得 $[T,a^*_{\alpha,n}]=-(n-1)a^*_{\alpha,n-1}$。结论：三处负号是同一约定的三个侧面，任一处改号都会使 (5.4-1) 与 (5.5-3) 整体差号。[来源：原书 p.142；Kac, \emph{Vertex Algebras for Beginners}, 2nd ed., AMS 1998, Ch.~1。]}
    \item 真空向量：$|0\rangle$；
    \item 平移算子：$T|0\rangle = 0$，$[T,a_{\alpha,n}] = -n\,a_{\alpha,n-1}$，$[T,a^*_{\alpha,n}] = -(n-1)a^*_{\alpha,n-1}$；
    \item 顶点算子：
\end{itemize}
\[
Y(a_{\alpha,-1}|0\rangle,z) = a_\alpha(z), \qquad Y(a^*_{\alpha,0}|0\rangle,z) = a^*_\alpha(z),
\]
\[
\begin{aligned}
Y(&a_{\alpha_1,n_1}\cdots a_{\alpha_k,n_k}a^*_{\beta_1,m_1}\cdots a^*_{\beta_l,m_l}|0\rangle,z)\\
&= \prod_{i=1}^{k}\frac{1}{(-n_i-1)!}\prod_{j=1}^{l}\frac{1}{(-m_j)!}\,
{:}\partial_z^{-n_1-1}a_{\alpha_1}(z)\cdots\partial_z^{-n_k-1}a_{\alpha_k}(z)\\
&\qquad\qquad\qquad\qquad\cdot\partial_z^{-m_1}a^*_{\beta_1}(z)\cdots\partial_z^{-m_l}a^*_{\beta_l}(z){:}.
\end{aligned}
\]

这里用到 2.2.5 节引入的正规序运算（书中记作冒号，形如两列竖点）。在一般情形，它是归纳定义的，因而并不算显式；但在当前情形可以给出显式的定义，我们下面回顾这一点。

% [Tutor Note: 分次约定极易抄错：书中写 deg a_{α,n} = deg a^*_{α,n} = −n，再配合平移算子的对易关系 [T,a_{α,n}] = −n a_{α,n−1} 与 [T,a^*_{α,n}] = −(n−1)a^*_{α,n−1}。文字层把这里的负号全部丢掉了，务必以图像为准。这套负号约定保证 T 提高次数 1（deg T = 1），是全章公式自洽的基准。]

我们把生成元 $a_{\alpha,n}$（$n\ge 0$）与 $a^*_{\alpha,m}$（$m>0$）称为\textbf{湮灭算子}，把 $a_{\alpha,n}$（$n<0$）与 $a^*_{\alpha,m}$（$m\le 0$）称为\textbf{产生算子}。$a_{\alpha,n}$ 与 $a^*_{\alpha,n}$ 中的单项式 $P$ 称为正规序的，如果 $P$ 中所有湮灭算子因子都排在所有产生算子因子的右边。给定任意单项式 $P$，定义正规序单项式 ${:}P{:}$ 为：把 $P$ 中所有湮灭算子因子移到右边、所有产生算子因子移到左边所得的单项式。例如
\[
{:}a_{\alpha,4}a^*_{\beta,-5}a^*_{\gamma,1}a^*_{\mu,-3}{:} = a^*_{\beta,-5}a^*_{\mu,-3}a_{\alpha,4}a^*_{\gamma,1}.
\]
注意，湮灭算子彼此对易，所以它们内部的排序方式无关紧要；产生算子同理。这说明 ${:}P{:}$ 由上述条件良定义。

给定两个单项式 $P$ 与 $Q$，它们的正规序乘积按定义就是正规序单项式 ${:}PQ{:}$。由线性性，顶点代数 $M_{\fg}$ 中任意多个顶点算子的正规序乘积，可对乘积每个 Fourier 系数中出现的每个单项式施用上述定义而得到。

现在令
\[
U(M_{\fg}) = (M_{\fg}\otimes \C((t)))/\operatorname{Im}(T\otimes 1 + \operatorname{Id}\otimes\partial_t)
\]
为如 3.2.1 节那样附着于 $M_{\fg}$ 的李代数。

正如 3.2.1 节所解释的，$U(M_{\fg})$ 可以看作 $M_{\fg}$ 中所有顶点算子的 Fourier 系数张成空间的完备化。此外，用与命题 3.2.1 相同的方法可以证明，映射
\[
U(M_{\fg}) \to \tilde{\cA}^{\fg},\qquad A\otimes f(z)\mapsto \Res_{z=0}Y(A,z)f(z)\,dz
\]
是李代数同态。

注意，$U(M_{\fg})$ 是李代数，但不是代数。例如它含有 Heisenberg 代数的生成元 $a_{\alpha,n}$、$a^*_{\alpha,n}$，却不含这些生成元的次数大于 1 的单项式。不过我们只需要 $U(M_{\fg})$ 上的李代数结构。\footnote{如同引理 3.2.4，由此可知 $\tilde{\cA}^{\fg}$ 同构于 $\widetilde{U}(M_{\fg})$，其中 $\widetilde{U}$ 是 3.2.3 节引入的函子；但我们不会用到这一事实。}

$\tilde{\cA}^{\fg}_0 = \operatorname{Fun} L\cU$ 中落在 $U(M_{\fg})$ 的像里的元素，通常称为 $L\cU$ 上的\textbf{局部泛函}。$\tilde{\cA}^{\fg}$ 中属于 $U(M_{\fg})$ 的元素，由 $a_\alpha(z)$、$a^*_\alpha(z)$ 及其导数的正规序多项式的 Fourier 系数（可能无穷的）线性组合给出，我们称之为 $\tilde{\cA}^{\fg}$ 的\textbf{局部元素}。

\subsection{$M_{\fg}$ 的更典范定义}

上面关于顶点代数 $M_{\fg}$ 的定义，用到了群 $N_+$ 上一个特定的坐标系 $y_\alpha$（$\alpha\in\Delta_+$）。若在 $N_+$ 上另选坐标系 $y'_\alpha$（$\alpha\in\Delta_+$），便得到另一个以 $a'_{\alpha,n}$、$a'^*_{\alpha,n}$ 为生成元的 Heisenberg 代数，以及顶点代数 $M'_{\fg}$。然而 $M'_{\fg}$ 与 $M_{\fg}$ 典范同构。特别地，用 $a_\alpha(z)$、$a^*_\alpha(z)$ 表达 $a'_\alpha(z)$、$a'^*_\alpha(z)$ 很容易：若 $y'_\alpha = F_\alpha(y_\beta)$，则
\[
a'_\alpha(z) \mapsto \sum_{\gamma\in\Delta_+}{:}\partial_{y_\gamma}F_\alpha(a^*_\beta(z))\,a_\gamma(z){:},\qquad
a'^*_\alpha(z)\mapsto F_\alpha(a^*_\beta(z)).
\]

注意，对坐标系的齐次性要求极大地限制函数 $F_\alpha$ 的形式：
\[
F_\alpha = c_\alpha y_\alpha + \sum_{\beta_1+\cdots+\beta_k=\alpha}c_{\beta_1,\dots,\beta_k}y_{\beta_1}\cdots y_{\beta_k},
\]
其中 $c_{\beta_1,\dots,\beta_k}\in\C$ 且 $c_\alpha\ne 0$。

也可以完全不参照 $N_+$ 上的坐标系来定义 $M_{\fg}$。具体地说，作为 $\fn_+((t))$-模，我们可以把 $M_{\fg}$ 等同于
\[
M_{\fg} = \operatorname{Ind}_{\fn_+[[t]]}^{\fn_+((t))}\operatorname{Fun}(N_+[[t]]) \simeq U(\fn_+\otimes t^{-1}\C[t^{-1}])\otimes \operatorname{Fun}(N_+[[t]]),
\]
其中 $\operatorname{Fun}(N_+[[t]])$ 是投射代数群 $N_+[[t]]$ 上的正则函数环，视为 $\fn_+[[t]]$-模。若选取 $N_+$ 上的坐标系
\[
\{y_\alpha\}_{\alpha\in\Delta_+},
\]
便得到 $N_+[[t]]$ 上的坐标系
\[
\{y_{\alpha,n}\}_{\alpha\in\Delta_+,\,n\ge 0}.
\]
于是 $u\otimes P(y_{\alpha,n})\in M_{\fg}$，其中
\[
u\in U(\fn_+\otimes t^{-1}\C[t^{-1}]) \quad\text{且}\quad P(y_{\alpha,n})\in\operatorname{Fun}(N_+[[t]]) = \C[y_{\alpha,n}]_{\alpha\in\Delta_+,\,n\ge 0},
\]
对应于 $M_{\fg}$ 先前描述中的 $u\cdot P(a^*_{\alpha,n})$。

在 $M_{\fg}$ 上定义顶点代数结构是直接的（见 [Feigin and Frenkel 1999] 第 2 节）。具体地说，$M_{\fg}$ 的真空向量是向量 $1\otimes 1\in M_{\fg}$。平移算子 $T$ 定义为算子 $-\partial_t$，它自然地作用在 $\operatorname{Fun}(N_+[[t]])$ 与 $\fn_+((t))$ 上，并保持 $\fn_+[[t]]$ 不变。接着，对有 $x\in\fn_+$ 的元素 $x_{-1}|0\rangle$，用公式
\[
Y(x_{-1}|0\rangle,z) = \sum_{n\in\mathbb{Z}}x_n z^{-n-1}
\]
定义相应的顶点算子，其中 $x_n = x\otimes t^n$，并把 $M_{\fg}$ 视为 $\fn_+((t))$ 的诱导表示来考虑其作用。我们还需要对 $P\in\operatorname{Fun}(N_+[[t]])$ 定义顶点算子
\[
Y(P|0\rangle,z) = \sum_{n\in\mathbb{Z}}P_{(n)}z^{-n-1}.
\]
相应的线性算子 $P_{(n)}$ 完全由它们在 $|0\rangle$ 上的作用、它们之间的对易性，以及它们与 $\fn_+((t))$ 的如下对易关系决定：
\[
P_{(n)}|0\rangle = 0,\quad n\ge 0;\qquad
P_{(n)}|0\rangle = \frac{1}{(-n-1)!}T^{-n-1}P|0\rangle;
\]
\[
[x_m, P_{(k)}] = \sum_{n\ge 0}\binom{m}{n}(x_n\cdot P)_{(m+k-n)}.
\]
利用重构定理 2.2.4 容易证明，这些公式在 $M_{\fg}$ 上定义了顶点代数结构。

% [Tutor Note: 这段“无坐标系”定义是全文最容易被略读的部分：M_g 作为 n_+((t))-模就是诱导模 Ind Fun(N_+[[t]])，顶点代数结构由真空 1⊗1、T = −∂_t 与生成元的顶点算子唯一确定，再用重构定理验证公理。注意 T = −∂_t（负号！）与 P_(n)|0⟩ = T^{−n−1}P|0⟩/(−n−1)! 中的阶乘符号，抄错会直接破坏分次。]

事实上，同样的定义在把 $N_+$ 换成任意代数群 $G$ 后仍然有效。在一般情形，自然要考虑圈代数 $\fg((t))$ 对应于 $\fg$ 上不变内积 $\kappa$ 的中心扩张 $\gghat_\kappa$，其中 $\kappa$ 按 5.3.3 节定义。于是有诱导模
\[
\operatorname{Ind}_{\fg[[t]]\oplus\C\mathbf{1}}^{\gghat_\kappa}\operatorname{Fun}(G[[t]]),
\]
其中中心元 $\mathbf{1}$ 在 $\operatorname{Fun}(G[[t]])$ 上按恒等映射作用。相应的顶点代数就是 $G$ 上的手性微分算子（chiral differential operators）代数，见 [Gorbounov et al. 2001; Arkhipov and Gaitsgory 2002]。如 [Gorbounov et al. 2001; Arkhipov and Gaitsgory 2002] 所证明，除 $\gghat_\kappa$ 在这个顶点代数上的自然（左）作用外，还有 $\gghat_{-\kappa-\kappa_{\fg}}$ 的另一个（右）作用，它与左作用对易。这里 $\kappa_{\fg}$ 是 $\fg$ 上的 Killing 型，由公式 $\kappa_{\fg}(x,y) = \Tr_{\fg}(\operatorname{ad}x\,\operatorname{ad}y)$ 定义。当 $\fg=\fn_+$ 时，不存在非零的不变内积（特别地 $\kappa_{\fn_+}=0$），于是得到 $\fn_+((t))$ 在 $M_{\fg}$ 上的一个对易右作用。下文将用到这个右作用（见 6.1.1 节）。\transnote{卡点：右作用的水平为什么是 $-\kappa-\kappa_{\fg}$ 而不是 $-\kappa$。背景：手性微分算子代数同时是 $\gghat_\kappa$ 的左模与 $\gghat_{-\kappa-\kappa_{\fg}}$ 的右模，两个作用彼此对易；右作用可看作“反向”的一份拷贝，而取对偶（contragredient）表示时水平按 $k\mapsto-k-2h^\vee$ 变化，写成不变内积即 $\kappa\mapsto-\kappa-\kappa_{\fg}$。结论：$\fg=\fn_+$ 时 $\kappa_{\fn_+}=0$（幂零李代数没有非零不变内积），右作用水平退化为 $0$，于是得到的是 $\fn_+((t))$ 本身（而非其中心扩张）的对易右作用；这正是 6.1.1 节定义 Wakimoto 模右结构时用到的情形。[来源：原书 p.145；Gorbounov--Malikov--Schechtman 2001、Arkhipov--Gaitsgory 2002（原书此处所引）。]}

% [Tutor Note: 右作用的水平是 −κ−κ_g（不是 −κ）：这是手性微分算子代数的对偶水平，第 6 章构造 Wakimoto 模的右作用时取的就是这个水平。对 g=n_+ 因无不变内积 κ_{n_+}=0，右作用退化为 n_+((t)) 本身的作用；这一点在 6.1.1 节定义 Wakimoto 模的右结构时会再次出现，是本书记号最容易混淆处之一。]

上面的公式实际上在
\[
\operatorname{Ind}_{\fn_+[[t]]}^{\fn_+((t))}\operatorname{Fun}(\cU[[t]])
\]
上定义了典范的顶点代数结构，它与 $N_+\simeq\cU$ 的等同选取无关。回忆一下：为了定义 $\cU$，只需固定 $G$ 的一个 Borel 子群 $B_+$。此时 $\cU$ 定义为旗簇中开的 $B_+$-轨道，因此自然是 $N_+$-torsor。要把 $\cU$ 与 $N_+$ 等同，需要选取 $\cU$ 中一点，即选取反 Borel 子群 $B_-$，或等价地选取 $B_+$ 的 Cartan 子群 $H = B_+\cap B_-$。但在上面的公式中，我们从未用到 $N_+$ 与 $\cU$ 的等同，只用到了 $\fn_+$ 在 $\cU$ 上的典范作用，而它决定了 $L\fn_+$ 在 $L\cU$ 上的典范作用。

若不固定等同 $N_+\simeq\cU$，则上述 $\fn_+$ 在 $N_+$ 上的右作用变成“扭变”李代数 $\fn_{+,\cU}$ 的作用，其中
\[
\fn_{+,\cU} = \cU\times_{N_+}\fn_+.
\]
有趣的是\footnote{感谢 D. Gaitsgory 指出这一点。}，与 $\fn_+$ 不同，这个扭变李代数 $\fn_{+,\cU}$ 有典范分解：它是相应于正根 $\alpha\in\Delta_+$ 的一维子空间 $\fn_{\alpha,\cU}$ 的直和。确实，对任意 $u\in\cU$，有等同 $\cU\simeq N_+$ 与 Cartan 子群 $H_u = B_+\cap B_u$，其中 $B_u$ 是 $u$ 的稳定化子。子空间 $\fn_{\alpha,\cU}$ 定义为 $\fn_+$ 关于 $H_u$ 的伴随作用的特征子空间。若另取一点 $u'\in\cU$，等同 $\cU\simeq N_+$ 与 $H_u$ 会改变，但特征子空间保持不变！\transnote{卡点：$H_u$ 随 $u$ 变化，为什么根空间 $\fn_{\alpha,\cU}$ 反而与 $u$ 无关。背景：任一 $u$ 给出 $\cU\simeq N_+$，而 $B_+$ 的任何 Cartan 子群都形如 $H_u=\operatorname{Ad}_n(H)$，$n\in N_+$。对 $x\in\fn_\alpha$，$\operatorname{Ad}_{n^{-1}}x=x+(\text{更高根项})$，故 $\operatorname{Ad}_{H_u}$ 在 $\fn_+$ 的根滤过上是三角的，对角特征标仍为 $\alpha$。结论：$\{\fn_{\alpha,\cU}\}$ 只依赖 $\fn_+$ 的 $B_+$-不变滤过，是滤过的分次分量，故为典范分解；这就是 $e_i^R$ 能“相差一个标量”良定义的原因。[来源：原书 p.146；译者按（上下文推断，无外部依据）。]}

因此，右作用 $\fn_{\alpha_i,\cU}$ 在 $\cU$ 上有典范（相差一个标量）的生成元 $e_i^R$。下文将用这些算子定义屏蔽算子；它们与 $B_+$ 中 Cartan 子群的选取无关，这一点意味着屏蔽算子的核不依赖于任何附加选取。

% [Tutor Note: 这里“特征子空间与 u 的选取无关”是屏蔽算子良定义的关键：扭变李代数 n_{+,U} 没有整体坐标系，却仍有典范的一维根空间分解，故可谈论 e_i^R（相差标量）。第 7 章定义屏蔽算子时“核与坐标选取无关”的结论，正是靠这条性质支撑的。]

\subsection{局部扩张}

对我们的目的而言，可以用相对较小的“局部部分” $U(M_{\fg})$ 替换掉非常大的拓扑代数 $\tilde{\cA}^{\fg}$。相应地，把 $\tilde{\cA}^{\fg}_0$ 与 $\tilde{\cA}^{\fg}_{\le1}$ 替换为它们的局部版本 $\cA^{\fg}_{0,\mathrm{loc}} = \tilde{\cA}^{\fg}_0\cap U(M_{\fg})$ 与 $\cA^{\fg}_{\le1,\mathrm{loc}} = \tilde{\cA}^{\fg}_{\le1}\cap U(M_{\fg})$。

更具体地描述 $\cA^{\fg}_{0,\mathrm{loc}}$ 与 $\cA^{\fg}_{\le1,\mathrm{loc}}$。空间 $\cA^{\fg}_{0,\mathrm{loc}}$ 由 $\partial_z^n a^*_\alpha(z)$（$n\ge 0$）的所有多项式的 Fourier 系数（拓扑地）张成。注意，由于 $a^*_{\alpha,n}$ 彼此对易，这些多项式自动是正规序的。空间 $\cA^{\fg}_{\le1,\mathrm{loc}}$ 由形如
\[
{:}P(a^*_\alpha(z),\partial_z a^*_\alpha(z),\dots)a_\beta(z){:}
\]
的场的 Fourier 系数张成（即 $P(a^*_\alpha(z),\partial_z a^*_\alpha(z),\dots)$ 与 $a_\beta(z)$ 的正规序乘积）。

这里用到了如下事实：所有形如
\[
{:}P(a^*_\alpha(z),\partial_z a^*_\alpha(z),\dots)\partial_z^m a_\beta(z){:},\qquad m>0,
\]
的场的 Fourier 系数，都可以表示成形如
\begin{equation}
{:}P(a^*_\alpha(z),\partial_z a^*_\alpha(z),\dots)a_\beta(z){:}
\tag{5.4-1}
\end{equation}
的场的 Fourier 系数的线性组合。

进一步，定义 $\operatorname{Vect} L\cU$ 的局部版本 $\cT^{\fg}_{\mathrm{loc}}$ 为如下子空间：它由形式幂级数
\begin{equation}
P(a^*_\alpha(z),\partial_z a^*_\alpha(z),\dots)a_\beta(z)
\tag{5.4-2}
\end{equation}
的 Fourier 系数的有限线性组合构成，其中 $a_\alpha(z)$ 与 $a^*_\alpha(z)$ 由公式 (5.3-4)、(5.3-5) 给出。

由于 $\cA^{\fg}_{\le1,\mathrm{loc}}$ 是 $\tilde{\cA}^{\fg}$ 的李子代数之交，它本身也是 $\tilde{\cA}^{\fg}$ 的李子代数。按构造，它在同态 $\tilde{\cA}^{\fg}_{\le1}\to\operatorname{Vect} L\cU$ 下的像等于 $\cT^{\fg}_{\mathrm{loc}}$。最后，所得满李代数同态 $\cA^{\fg}_{\le1,\mathrm{loc}}\to\cT^{\fg}_{\mathrm{loc}}$ 的核等于 $\cA^{\fg}_{0,\mathrm{loc}}$。于是得到：扩张 (5.3-7) 限制为“局部”扩张
\begin{equation}
0 \to \cA^{\fg}_{0,\mathrm{loc}} \to \cA^{\fg}_{\le1,\mathrm{loc}} \to \cT^{\fg}_{\mathrm{loc}} \to 0.
\tag{5.4-3}
\end{equation}
下文将看到这个序列不可分裂。相应的二上循环将在引理 5.5.2 中用 Wick 公式显式计算（它来自相应顶点算子的“双缩并”）。

按 5.3.2 节，$L\fg$ 在 $\operatorname{Vect} L\cU$ 中的像落在 $\cT^{\fg}_{\mathrm{loc}}$ 里。我们将证明，同态 $L\fg\to\cT^{\fg}_{\mathrm{loc}}$ 可以提升为同态 $\gghat_\kappa\to\cA^{\fg}_{\le1,\mathrm{loc}}$，其中 $\gghat_\kappa$ 是 5.3.3 节定义的 $L\fg$ 的中心扩张。

% [Tutor Note: (5.4-1) 与 (5.4-2) 的差别只有一层正规序，但正是这一层正规序给出了向量空间分裂 ι；局部扩张 (5.4-3) 就是本节的全部内容。注意 (5.4-2) 是“向量场”的生成函数，本身不是场（无正规序），所以不能直接对其用 OPE，必须先经 ι 提升到 A_{≤1,loc}。]

\section{二上循环的计算}

\begin{keypoints}
    \item \textbf{核心思想：}用正规序作为向量空间分裂，把局部扩张的不可分裂性量化为一个二上循环，再用 Wick 公式显式计算它。
    \item \textbf{数学工具：}算子乘积展开（OPE）、Wick 公式、双缩并、Chevalley 复形与 $H^2(\fg,\fh)$。
    \item \textbf{关键结论：}式 (5.5-3) 给出二上循环 $\omega$ 的显式公式；对 Cartan 型元素算出 $\omega(\mathbf{h}_n,\mathbf{h}_m) = -n\delta_{n,-m}$，从而 (5.3-7) 与 (5.4-3) 均不可分裂。
    \item \textbf{记号提醒：}$\omega$ 专指局部扩张 (5.4-3) 的二上循环；5.5 节一般的二上循环记作 $\sigma$，而 5.5.5 节又用 $\sigma$ 表示内积 $\kappa_c$ 对应的上同调类。
\end{keypoints}

回忆：李代数正合列
\[
0\to \mathfrak{h}\to \tilde{\mathfrak{g}}\to \fg\to 0,
\]
其中 $\mathfrak{h}$ 是阿贝尔理想，并带有给定的 $\fg$-模结构，会给出 $\fg$ 的一个取值于 $\mathfrak{h}$ 的二上循环。构造如下：选取该序列作为向量空间的一个分裂 $\iota: \fg\to\tilde{\mathfrak{g}}$，并定义 $\sigma: \bigwedge^2\fg\to\mathfrak{h}$ 为
\[
\sigma(a,b) = \iota([a,b]) - [\iota(a),\iota(b)].
\]
可以验证，$\sigma$ 是 $\fg$ 的取值于 $\mathfrak{h}$ 的 Chevalley 复形中的一个二上循环；改变分裂 $\iota$ 相当于把 $\sigma$ 改变一个上边缘。

反过来，给定线性泛函 $\sigma:\bigwedge^2\fg\to\mathfrak{h}$，可以在直和 $\fg\oplus\mathfrak{h}$ 上构造李代数结构：$\mathfrak{h}$ 中任意两元素的李括号为 0，且对一切 $X\in\fg$、$h\in\mathfrak{h}$ 有 $[X,h] = X\cdot h$，此外
\[
[X,Y] = [X,Y]_{\fg} + \sigma(X,Y),\qquad X,Y\in\fg.
\]
这些公式在 $\fg\oplus\mathfrak{h}$ 上定义李代数结构，当且仅当 $\sigma$ 是 $\fg$ 的标准 Chevalley 复形中取值于 $\mathfrak{h}$ 的二上循环。因此得到：$\fg$ 通过 $\mathfrak{h}$ 的扩张的同构类集合与上同调群 $H^2(\fg,\mathfrak{h})$ 一一对应。

% [Tutor Note: 这里把“扩张 ↔ 二上循环”的对应写清了，但符号约定必须记牢：σ(a,b) = ι([a,b]) − [ι(a),ι(b)]，即“先括号再提升”减“先提升再括号”。若把 ι 换成 −ι 或把次序反过来，全节所有二上循环都会差一个符号，5.5.3 与 5.6.4 的结论也会随之翻转，核对时以此处为基准。]

\subsection{Wick 公式}

考虑扩张 (5.4-3)。正规序运算给出该扩张作为向量空间的一个分裂 $\iota$：具体地说，$\iota$ 把级数 (5.4-2) 的第 $n$ 个 Fourier 系数映到级数 (5.4-1) 的第 $n$ 个 Fourier 系数。为计算相应的二上循环，必须学会计算形如 (5.4-1) 与 (5.4-2) 的生成函数的 Fourier 系数的对易子。这些可由相应顶点算子的算子乘积展开（OPE）算出。下面解释如何用 Wick 公式计算顶点算子的 OPE。

为陈述 Wick 公式，需要引入两个场的缩并概念。为简化记号，假设 $\fg=\fsl_2$，并略去 $a_{\alpha,n}$ 与 $a^*_{\alpha,n}$ 中的指标 $\alpha$；一般情形处理方式相同。

由对易关系得到如下 OPE：
\[
a(z)a^*(w) = \frac{1}{z-w} + {:}a(z)a^*(w){:},\qquad
a^*(z)a(w) = -\frac{1}{z-w} + {:}a^*(z)a(w){:}.
\]
现在把它们看作形式幂级数之间的恒等式，其中 $1/(z-w)$ 理解为按 $w/z$ 的非负幂展开。求导若干次，得到
\begin{equation}
\partial_z^n a(z)\,\partial_w^m a^*(w) = (-1)^n\frac{(n+m)!}{(z-w)^{n+m+1}} + {:}\partial_z^n a(z)\,\partial_w^m a^*(w){:},
\tag{5.5-1}
\end{equation}
\begin{equation}
\partial_z^m a^*(z)\,\partial_w^n a(w) = (-1)^{m+1}\frac{(n+m)!}{(z-w)^{n+m+1}} + {:}\partial_z^m a^*(z)\,\partial_w^n a(w){:}
\tag{5.5-2}
\end{equation}
（这里同样把 $1/(z-w)^n$ 理解为按 $w/z$ 的非负幂展开）。\transnote{卡点：$a$ 与 $a^*$ 互为共轭场，为什么 $a(z)a^*(w)$ 的 OPE 是 $+1/(z-w)$ 而 $a^*(z)a(w)$ 的是 $-1/(z-w)$。背景：两个场的模指标整体错开一位（$a$ 的场取 $z^{-n-1}$，$a^*$ 的场取 $z^{-n}$），而它们的模满足 $[a_n,a^*_m]=\delta_{n+m,0}$；按相反次序相乘时，参与缩并的模指标范围随之错开，单缩并的常数项恰好变号。结论：这个负号不是笔误，它是 (5.5-2) 中 $(-1)^{m+1}$ 以及 (5.5-3) 全部双缩并项定号的来源；若把第二式改成 $+1/(z-w)$，后面 (5.5-3) 的每个双缩并项都会整体变号。[来源：原书 p.148；(5.5-1)、(5.5-2) 即对这两个 OPE 逐次求导所得。]}

% [Tutor Note: 符号高危处：a^*(z)a(w) 的 OPE 系数是 −1/(z−w)（负号！），(5.5-2) 前的 (−1)^{m+1} 就源于此；若误写成 +1/(z−w)，则 (5.5-2) 变成 (−1)^m，随后 (5.5-3) 的所有双缩并项与 5.5.3 的 −nδ_{n,−m} 都会错。文字层把 (5.5-2) 的负号丢掉了，必须看图像。]

设有两个关于 $a(z)$、$a^*(z)$ 及其导数的正规序单项式，记作 $P(z)$ 与 $Q(z)$。$P(z)$ 与 $Q(w)$ 之间的一个\textbf{单配对}，按定义是指：$P(z)$ 中出现的 $\partial_z^n a(z)$ 与 $Q(w)$ 中出现的 $\partial_w^m a^*(w)$ 构成的配对 $(\partial_z^n a(z),\partial_w^m a^*(w))$；或者 $P(z)$ 中出现的 $\partial_z^m a^*(z)$ 与 $Q(w)$ 中出现的 $\partial_w^n a(w)$ 构成的配对 $(\partial_z^m a^*(z),\partial_w^n a(w))$。我们分别给它配上函数
\[
(-1)^n\frac{(n+m)!}{(z-w)^{n+m+1}}
\qquad\text{与}\qquad
(-1)^{m+1}\frac{(n+m)!}{(z-w)^{n+m+1}}.
\]
\textbf{多重配对} $B$ 按定义是若干互不相交的单配对的并。给它配上函数 $f_B(z,w)$，即 $B$ 中各单配对对应函数之积。

注意，单项式 $P(z)$ 与 $Q(z)$ 可能含有同型的多个配对。例如单项式 ${:}a^*(z)^2\partial_z a(z){:}$ 与 ${:}a(w)\partial_w^2 a^*(w){:}$ 有两个不同的 $(a^*(z),a(w))$ 型配对，相应的函数为 $-1/(z-w)$；此时我们说这个配对的重数为 2。注意这两个单项式还有唯一的配对 $(\partial_z a(z),\partial_w^2 a^*(w))$，相应函数为 $-6/(z-w)^4$。

给定 $P(z)$ 与 $Q(w)$ 之间的多重配对 $B$，把 $(P(z)Q(w))_B$ 定义为 $P(z)$ 与 $Q(w)$ 中不属于该配对的所有因子之积（若没有剩余因子，则令 $(P(z)Q(w))_B = 1$）。$P(z)Q(w)$ 关于配对 $B$ 的\textbf{缩并}记作 ${:}P(z)Q(w){:}_B$，按定义是正规序形式幂级数 ${:}(P(z)Q(w))_B{:}$ 乘以函数 $f_B(z,w)$。当 $B$ 为空集时，补充约定
\[
{:}P(z)Q(w){:}_\varnothing = {:}P(z)Q(w){:}.
\]

现在可以陈述\textbf{Wick 公式}了，它给出两个任意正规序单项式顶点算子的 OPE。该公式的证明直接，留给读者。

\begin{lemma}
设 $P(z)$ 与 $Q(w)$ 是如上两个单项式。则乘积 $P(z)Q(w)$ 等于 $\sum_B {:}P(z)Q(w){:}_B$，求和遍历 $P$ 与 $Q$ 之间的一切配对 $B$（包括空配对），并按重数计。
\end{lemma}

下面是一个例子：
\[
\begin{aligned}
{:}a^*(z)^2\partial_z a(z){:}\,{:}a(w)\partial_z^2 a^*(w){:}
={}&{:}a^*(z)^2\partial_z a(z)a(w)\partial_z^2 a^*(w){:}\\
&-\frac{2}{z-w}{:}a^*(z)\partial_z a(z)\partial_z^2 a^*(w){:}
-\frac{6}{(z-w)^4}{:}a^*(z)^2a(w){:}\\
&+\frac{12}{(z-w)^5}a^*(z){:}.
\end{aligned}
\]

\subsection{双缩并}

现在可以计算二上循环了。为此，需要对形如
\[
{:}R(a^*(z),\partial_z a^*(z),\dots)a(z){:}
\]
的场应用 Wick 公式，这些场的 Fourier 系数在我们的分裂 $\iota$ 下张成 $\cT_{\mathrm{loc}}$ 在 $\cA_{\le1,\mathrm{loc}}$ 中的原像。两个这种形式的场只可能有单配对或双配对，所以它们的 OPE 可以写得相当显式。

上述形式的场可以写成 $Y(P(a^*_n)a_{-1},z)$（简记为 $Y(Pa_{-1},z)$），其中 $P$ 是 $a^*_n$（$n\le 0$）的多项式（回忆 $a^*_n$ 对应 $\partial_z^{-n}a^*(z)/(-n)!$）。\footnote{下文中，我们滥用记号，把严格说来应记作 $Y[A,z]$ 的级数写成 $Y(A,z)$。理由是同态 $U(M_{\fg})\to\operatorname{End}M_{\fg}$、$A_{[n]}\mapsto A_{(n)}$ 是单射，因此考虑 $U(M_{\fg})$ 在 $\operatorname{End}M_{\fg}$ 中的像不会丢失任何信息。} 应用 Wick 公式，得到

\begin{lemma}
\begin{align*}
Y(Pa_{-1},z)Y(Qa_{-1},w)
&= {:}Y(Pa_{-1},z)Y(Qa_{-1},w){:}\\
&\quad+\sum_{n\ge 0}\frac{1}{(z-w)^{n+1}}{:}Y(P,z)Y\Big(\frac{\partial Q}{\partial a^*_{-n}}a_{-1},w\Big){:}\\
&\quad-\sum_{n\ge 0}\frac{1}{(z-w)^{n+1}}{:}Y\Big(\frac{\partial P}{\partial a^*_{-n}}a_{-1},z\Big)Y(Q,w){:}\\
&\quad-\sum_{n,m\ge 0}\frac{1}{(z-w)^{n+m+2}}Y\Big(\frac{\partial P}{\partial a^*_{-n}},z\Big)Y\Big(\frac{\partial Q}{\partial a^*_{-m}},w\Big).
\end{align*}
\end{lemma}

注意，最后一项不需要加正规序。\transnote{卡点：下面 (5.5-3) 的求值次序，以及 $|_{z=w}$ 究竟作用在哪一步。背景：$\partial_z^{n+m+1}$ 只作用于第一个因子 $Y(\partial P/\partial a^*_{-n},z)$，第二个因子 $Y(\partial Q/\partial a^*_{-m},w)$ 在此步视为常数；求导完毕才令 $z=w$；最后把结果看作 $w$ 的形式幂级数并取 $\operatorname{Res}_{w=0}$（即 $w^{-1}$ 的系数）。结论：先求导、后代入、再取留数，三步不可交换。这一项正是双缩并，没有正规序保护，因而是整个二上循环的唯一来源。[来源：原书 p.150（(5.5-3) 即引理 5.5.2 最后一项的 Fourier 系数形式）。]}

利用这个公式与对易关系 (3.2-2)，现在容易得到场 $Y(Pa_{-1},z)$ 与 $Y(Qa_{-1},w)$ 的 Fourier 系数的对易子。

引理 5.5.2 右端前两项对应 $Y(Pa_{-1},z)$ 与 $Y(Qa_{-1},w)$ 之间的单缩并。这些项在 Fourier 系数对易子中贡献的那部分，恰好等于在 $\cT_{\mathrm{loc}}$ 中算出的相应向量场的对易子。于是我们看到，$\cA_{\le1,\mathrm{loc}}$ 与 $\cT_{\mathrm{loc}}$ 中对易子的差异（由二上循环度量）来自引理 5.5.2 公式的最后一项，它来自 $Y(Pa_{-1},z)$ 与 $Y(Qa_{-1},w)$ 之间的\textbf{双缩并}。

显式地，二上循环由如下公式给出
\begin{equation}
\begin{aligned}
\omega(&(Pa_{-1})_{[k]},(Qa_{-1})_{[s]})\\
&=-\sum_{n,m\ge 0}\Res_{w=0}\frac{1}{(n+m+1)!}\,
\partial_z^{n+m+1}Y\Big(\frac{\partial P}{\partial a^*_{-n}},z\Big)Y\Big(\frac{\partial Q}{\partial a^*_{-m}},w\Big)\\
&\qquad\qquad\qquad\qquad\qquad\qquad\cdot z^kw^s\Big|_{z=w}\,dw.
\end{aligned}
\tag{5.5-3}
\end{equation}

\subsection{扩张不可分裂}

例如，对由公式
\[
Y(a^*_0a_{-1},z) = {:}a^*(z)a(z){:} = \sum_{n\in\mathbb{Z}}\mathbf{h}_n z^{-n-1}
\]
定义的元 $\mathbf{h}_n$ 计算上循环，于是
\[
\mathbf{h}_n = \sum_{k\in\mathbb{Z}}{:}a^*_{-k}a_{n+k}{:}.
\]
按公式 (5.5-3)，有
\[
\omega(\mathbf{h}_n,\mathbf{h}_m) = -\Res_{w=0}n\,w^{n+m-1}dw = -n\delta_{n,-m}.
\]

% [Tutor Note: (5.5-3) 的求值记号要当心：∂_z^{n+m+1} 只作用在 Y(∂P/∂a^*_{−n}, z) 上，随后才令 z=w；文字层把 |_{z=w} 的位置弄乱了。另外这一行 ω(h_n,h_m) = −Res n w^{n+m−1} dw = −nδ_{n,−m} 的负号与 δ_{n,−m}（不是 δ_{n,m}）都是关键，抄成 +nδ_{n,m} 会使 5.5.3 的不可分裂论证失效。]

由于单缩并项相互抵消，得到 $\cA_{\le1,\mathrm{loc}}$ 中的如下对易关系：
\[
[\mathbf{h}_n,\mathbf{h}_m] = -n\delta_{n,-m}.
\]
另一方面，$\mathbf{h}_n$ 在 $\cT_{\mathrm{loc}}$ 中的像是向量场
\[
\bar{\mathbf{h}}_n = \sum_{k\in\mathbb{Z}}y_k\frac{\partial}{\partial y_{n+k}}.
\]
这些向量场彼此对易。

若序列 (5.3-7)（或 (5.4-3)）分裂，就能找到“修正项” $f_m\in\operatorname{Fun}\C((t))$（或 $\cA_{0,\mathrm{loc}}$ 中的 $f_m$），使得提升 $\mathbf{h}_n\mapsto\mathbf{h}_n+f_n$ 保持李括号，即
\[
[\mathbf{h}_n+f_n,\mathbf{h}_m+f_m] = 0,\qquad \forall n,m\in\mathbb{Z}.
\]
但这等价于公式
\[
\bar{\mathbf{h}}_n\cdot f_m - \bar{\mathbf{h}}_m\cdot f_n = -n\delta_{n,-m}.
\]
% [Tutor Note: 此处符号有疑：由上一行 [h_n,h_m] = −nδ_{n,−m} 与 [h_n,f_m] = h̄_n·f_m，令 [h_n+f_n,h_m+f_m]=0 应得 h̄_n·f_m − h̄_m·f_n = +nδ_{n,−m}，与印出的 −nδ_{n,−m} 差一个符号。译文保留印出的形式，请以本批注为准。]
然而，由于 $\bar{\mathbf{h}}_n$ 是线性向量场（即关于 $L\cU$ 上的坐标 $y_k$ 线性），无论怎样选取 $f_m$（$n\in\mathbb{Z}$），上式左端都不可能成为非零常数。\transnote{卡点：上式与上一段的 $[\mathbf{h}_n,\mathbf{h}_m]=-n\delta_{n,-m}$ 差一个符号。背景：把 $[\mathbf{h}_n+f_n,\mathbf{h}_m+f_m]=0$ 展开（$f$ 之间对易，$[\mathbf{h}_n,f_m]=\bar{\mathbf{h}}_n\cdot f_m$），得到的是 $\bar{\mathbf{h}}_n\cdot f_m-\bar{\mathbf{h}}_m\cdot f_n=+n\delta_{n,-m}$。原书 p.151 同页先后印出这两个相差一号的式子，属原书笔误；此外紧随其后的括号里“$n\in\Z$”应为“$m\in\Z$”（$f_m$ 的指标）。结论：译文照原书保留印出的负号，正确读法是左端等于 $+n\delta_{n,-m}$；下一句“左端不可能是非零常数”的论证不受影响（两种符号下都不可能）。[来源：原书 p.151（PDF p.164）同页两式。]}

于是可得：序列 (5.3-7) 与 (5.4-3) 作为李代数的正合序列都不可分裂。

\subsection{上同调的回顾}

因此，不能指望像有限维情形那样把同态 $\widehat{\rho}: L\fg\to\cT_{\mathrm{loc}}$ 提升为同态 $L\fg\to\tilde{\cA}^{\fg}_{\le1,\mathrm{loc}}$。次优的目标是把它提升为从 $L\fg$ 的中心扩张 $\gghat_\kappa$ 到 $\tilde{\cA}^{\fg}_{\le1,\mathrm{loc}}$ 的同态。下面看看这样的提升存在性对应什么上同调条件。

仍设
\[
0\to\mathfrak{h}\to\tilde{\mathfrak{l}}\to\mathfrak{l}\to 0
\]
是李代数扩张，其中 $\mathfrak{h}$ 是阿贝尔李子代数，且在 $\tilde{\mathfrak{l}}$ 中是理想。取该序列作为向量空间的分裂 $\iota$，如 5.5 节那样定义 $\mathfrak{l}$ 的取值于 $\mathfrak{h}$ 的二上循环。设给定某个李代数 $\fg$ 的李代数同态 $\alpha:\fg\to\mathfrak{l}$。把这个二上循环沿 $\alpha$ 拉回，就得到 $\fg$ 的取值于 $\mathfrak{h}$ 的二上循环。

另一方面，给定 $\fg$-模的同态 $\mathfrak{h}'\xrightarrow{i}\mathfrak{h}$，便得到 $\fg$ 的取值于 $\mathfrak{h}$ 与 $\mathfrak{h}'$ 的二上循环空间之间的映射 $i_*$。相应的上同调群映射 $H^2(\fg,\mathfrak{h}')\to H^2(\fg,\mathfrak{h})$ 也记作 $i_*$。

\begin{lemma}
设给定 $\fg$ 的一个取值于 $\mathfrak{h}'$ 的二上循环 $\sigma$，使得 $i_*(\sigma)$ 与 $\omega$ 的上同调类在 $H^2(\fg,\mathfrak{h}')$ 中相等。按上述方式，用 $\sigma$ 定义 $\fg$ 通过 $\mathfrak{h}'$ 的扩张 $\tilde{\fg}$。则映射 $\fg\to\mathfrak{l}$ 可以扩充为交换图表
\begin{equation}
\begin{CD}
0 @>>> \fh @>>> \tilde{\mathfrak{l}} @>>> \mathfrak{l} @>>> 0\\
@. @AAA @AAA @AAA @.\\
0 @>>> \fh' @>>> \tilde{\fg} @>>> \fg @>>> 0
\end{CD}
\tag{5.5-4}
\end{equation}
的映射。此外，这类图表的同构类构成 $H^1(\fg,\mathfrak{h})$ 上的一个 torsor。\transnote{卡点：陈述中“在 $H^2(\fg,\mathfrak{h}')$ 中相等”一句，$\mathfrak{h}'$ 疑为 $\mathfrak{h}$ 之误。背景：$i_*$ 由 $\mathfrak{h}'\xrightarrow{i}\mathfrak{h}$ 诱导，把 $H^2(\fg,\mathfrak{h}')$ 送到 $H^2(\fg,\mathfrak{h})$；$\sigma$ 取值于 $\mathfrak{h}'$，$\omega$ 取值于 $\mathfrak{h}$，二者只能在 $H^2(\fg,\mathfrak{h})$ 中比较，证明第一句所写的 $i_*(\sigma)+d\gamma=\omega$（$\gamma:\fg\to\mathfrak{h}$）也要求如此。结论：译文照原书保留 $\mathfrak{h}'$，读者应按 $H^2(\fg,\mathfrak{h})$ 理解；引理的内容与后面 torsor 的结论不受影响。[来源：原书 p.152 的陈述与证明首句（PDF p.165）。]}
\end{lemma}
% [Tutor Note: 引理 5.5.3 的陈述里“在 H²(g,h′) 中相等”疑为“在 H²(g,h) 中相等”之误：i_*(σ) 已取值于 h，而证明中写 i_*(σ)+dγ=ω 也要求二者同处 H²(g,h)。译文保留原文，此处登记。]

\begin{proof}
若 $i_*(\sigma)$ 与 $\omega$ 的上同调类相同，则 $i_*(\sigma)+d\gamma = \omega$，其中 $\gamma$ 是一上链，即线性泛函 $\fg\to\mathfrak{h}$。如下定义线性映射 $\beta:\tilde{\fg}\to\tilde{\mathfrak{l}}$：按定义有向量空间分裂 $\tilde{\fg} = \fg\oplus\mathfrak{h}'$。对一切 $X\in\fg$ 令 $\beta(X) = \iota(\alpha(X))+\gamma(X)$，对一切 $h\in\mathfrak{h}$ 令 $\beta(h) = i(h)$。则上述上循环之相等蕴含 $\beta$ 是李代数同态，且使图表 (5.5-4) 交换。然而 $\gamma$ 的选取不唯一，可以给它加上任意一上循环 $\gamma'$。但对应 $\gamma$ 与 $\gamma+\gamma'$（其中 $\gamma'$ 是上边缘）的同态给出同构的图表。由此可知这类图表的同构类构成 $H^1(\fg,\mathfrak{h})$ 上的一个 torsor。
\end{proof}

\subsection{两个上循环}

把对应扩张 (5.4-3) 的、$\cT_{\mathrm{loc}}$ 的取值于 $\cA^{\fg}_{0,\mathrm{loc}}$ 的二上循环 $\omega$ 限制到 $L\fg\subset\cT_{\mathrm{loc}}$，得到 $L\fg$ 的取值于 $\cA^{\fg}_{0,\mathrm{loc}}$ 的二上循环，仍记作 $\omega$。$L\fg$-模 $\cA^{\fg}_{0,\mathrm{loc}}$ 含有平凡子表示 $\C$，即 $|0\rangle_{[-1]}$ 张成的一维空间（把它看作 $L\cU$ 上的常函数）；包含映射 $\C\xrightarrow{i}\cA^{\fg}_{0,\mathrm{loc}}$ 诱导出相应二上循环空间与上同调群之间的映射 $i_*: H^2(L\fg,\C)\to H^2(L\fg,\cA^{\fg}_{0,\mathrm{loc}})$。

正如 5.3.3 节所讨论的，上同调群 $H^2(L\fg,\C)$ 是一维的，同构于 $\fg$ 上不变内积构成的空间。把对应内积 $\kappa_c = -\tfrac12\kappa_{\fg}$ 的上同调类记作 $\sigma$，其中 $\kappa_{\fg}$ 表示 $\fg$ 上的 Killing 型。于是按定义
\[
\kappa_c(x,y) = -\frac{1}{2}\Tr(\operatorname{ad}x\,\operatorname{ad}y).
\]
此时上循环 $\sigma$ 由公式
\begin{equation}
\sigma(J^a_n,J^b_m) = n\delta_{n,-m}\kappa_c(J^a,J^b)
\tag{5.5-5}
\end{equation}
给出（见 5.3.3 节）。\transnote{卡点：$\kappa_c=-\tfrac12\kappa_{\fg}$ 究竟是多“特殊”的水平。背景：若把 $\kappa_{\fg}$ 写成 $2h^\vee\kappa_0$（$\kappa_0$ 为 $(\theta,\theta)=2$ 的规范化不变内积），则 $\kappa_c=-h^\vee\kappa_0$，即仿射 Kac--Moody 代数的临界水平 $k=-h^\vee$；此处 Sugawara 构造退化，顶点代数的中心由真空模的完备化给出，这正是第 8--9 章的主线。结论：负号与 $\tfrac12$ 都不是笔误，它们决定 (5.5-5) 与 5.6 节全部比较的符号——$\kappa_c$ 为负，正是 (5.6-11) 必须带负号的原因（见原书 p.161 的推导）。[来源：原书 p.152；Kac, \emph{Infinite-Dimensional Lie Algebras}, 3rd ed., Cambridge Univ. Press 1990（临界水平与 Sugawara 构造）。]}

% [Tutor Note: 临界内积的约定必须记牢：κ_c = −(1/2)κ_g，κ_g(x,y) = Tr(ad x ad y)。等价于 k = −h^∨（在 (θ,θ)=2 的规范化下）。本节的 −1/2 与 (5.5-5) 的 nδ_{n,−m} 一起决定了 5.6 节所有比较的符号；若把 κ_c 写成 +(1/2)κ_g，则 (5.6-11) 反而“对”了，但那说明的是约定被改动，不是原书自洽。]

\section{上同调类的比较}

\begin{keypoints}
    \item \textbf{核心思想：}$H^2(L\fg,\cA^{\fg}_{0,\mathrm{loc}})$ 的 Chevalley 复形过大，先把 $\omega$ 与 $i_*(\sigma)$ 放进一个可管理的小复形 $\cC^\bullet_{\mathrm{loc}}$，再证明它们限制到 $L\fh$ 上相等，从而同调。
    \item \textbf{数学工具：}Clifford 顶点超代数 $\bigwedge_{\fg}$、局部 Chevalley 复形、Shapiro 引理、Serre--Hochschild 谱序列。
    \item \textbf{关键结论：}定理 5.6.8：$\omega$ 与 $i_*(\sigma)$ 在 $H^2_{\mathrm{loc}}$ 中代表同一上同调类，故存在自由场实现同态 $\gghat_{\kappa_c}\to\cA^{\fg}_{\le1,\mathrm{loc}}$、$\mathbf{1}\mapsto 1$。
    \item \textbf{记号提醒：}$L_+\fg = \fg[[t]]$；$L_+\fh$ 的上同调被识别为 $\bigwedge^\bullet(L_+\fh)^*$，其中 $\fh$ 是 Cartan 子代数。
\end{keypoints}

下一节将证明 $i_*(\sigma)$ 与 $\omega$ 的上同调类相等。届时引理 5.5.3 将蕴含：存在一族李代数同态 $\gghat_{\kappa_c}\to\cA^{\fg}_{\le1,\mathrm{loc}}$，使得 $\mathbf{1}\mapsto 1$。

可惜，计算 $H^2(L\fg,\cA^{\fg}_{0,\mathrm{loc}})$ 的 Chevalley 复形大得无法处理。所以第一步先在下一小节证明：$\omega$ 与 $\sigma$ 其实都属于一个小得多的子复形，在那里比较它们要容易得多。

\subsection{Clifford 代数}

选取 $\fg$ 的一组基 $\{J^a\}_{a=1,\dots,\dim\fg}$，并令 $J^a_n = J^a\otimes t^n$。引入以
\[
\psi_{a,n},\ \psi^*_{a,m},\quad a=1,\dots,\dim\fg;\quad m,n\in\mathbb{Z}
\]
为生成元的 Clifford 代数，其反对易关系为
\[
[\psi_{a,n},\psi_{b,n}]_+ = [\psi^*_{a,n},\psi^*_{b,m}]_+ = 0,\qquad
[\psi_{a,n},\psi^*_{b,m}]_+ = \delta_{a,b}\delta_{n,-m}.
\]
设 $\bigwedge_{\fg}$ 是该 Clifford 代数由向量 $|0\rangle$ 生成的模，其中
\[
\psi_{a,n}|0\rangle = 0,\quad n>0;\qquad \psi^*_{a,n}|0\rangle = 0,\quad n\ge 0.
\]
于是 $\bigwedge_{\fg}$ 带有如下 $\mathbb{Z}_+$-分次顶点超代数结构（见 [Frenkel and Ben-Zvi 2004] 第 15.1.1 节）：

\begin{itemize}
    \item $\mathbb{Z}_+$-分次：$\deg\psi_{a,n} = \deg\psi^*_{a,n} = -n$，$\deg|0\rangle = 0$；
    \item 真空向量：$|0\rangle$；
    \item 平移算子：$T|0\rangle = 0$，$[T,\psi_{a,n}] = -n\psi_{a,n-1}$，$[T,\psi^*_{a,n}] = -(n-1)\psi^*_{a,n-1}$；
    \item 顶点算子：
\end{itemize}
\[
Y(\psi_{a,-1}|0\rangle,z) = \psi_a(z) = \sum_{n\in\mathbb{Z}}\psi_{a,n}z^{-n-1},\qquad
Y(\psi^*_{a,0}|0\rangle,z) = \psi^*_a(z) = \sum_{n\in\mathbb{Z}}\psi^*_{a,n}z^{-n},
\]
\[
\begin{aligned}
Y(&\psi_{a_1,n_1}\cdots\psi_{a_k,n_k}\psi^*_{b_1,m_1}\cdots\psi^*_{b_l,m_l}|0\rangle,z)\\
&= \prod_{i=1}^{k}\frac{1}{(-n_i-1)!}\prod_{j=1}^{l}\frac{1}{(-m_j)!}\,
{:}\partial_z^{-n_1-1}\psi_{a_1}(z)\cdots\partial_z^{-n_k-1}\psi_{a_k}(z)\\
&\qquad\qquad\qquad\qquad\qquad\cdot\partial_z^{-m_1}\psi^*_{b_1}(z)\cdots\partial_z^{-m_l}\psi^*_{b_l}(z){:}.
\end{aligned}
\]

% [Tutor Note: Clifford 侧的约定与 Heisenberg 侧平行，但有一处极易错：ψ^* 的顶点算子幂次是 z^{−n}（没有 −1），而 ψ 的是 z^{−n−1}；反对易关系配对的是 δ_{n,−m} 而不是 δ_{n,m}。抄错这两处会让 (5.6-4) 的 i_*(σ) 表达式对不上。]

两个顶点超代数的张量积自然仍是顶点超代数（见 [Frenkel and Ben-Zvi 2004] 引理 1.3.6），因此 $M_{\fg}\otimes\bigwedge_{\fg}$ 是顶点超代数。\transnote{卡点：Clifford 侧为什么要用两套不同的幂次——$\psi$ 的场取 $z^{-n-1}$、$\psi^*$ 的场取 $z^{-n}$，且反对易子是 $\delta_{n,-m}$ 而非 $\delta_{n,m}$。背景：这与 Heisenberg 侧完全平行，$\psi$ 与 $\psi^*$ 是一对共轭场，模指标整体错开一位；取 $z^{-n}$ 的那个场对应的是 $\psi^*_{a,0}|0\rangle$（而非 $\psi^*_{a,-1}|0\rangle$）。结论：这套约定保证 $Y(T\psi,z)=\partial_zY(\psi,z)$ 与 $\bigwedge_{\fg}$ 的 $\Z_+$-分次相容；若把 $\delta_{n,-m}$ 写成 $\delta_{n,m}$，(5.6-4) 中的 $\psi^*_{a,-n}\psi^*_{b,n}$ 就无法配对。[来源：原书 p.153；Frenkel--Ben-Zvi, \emph{Vertex Algebras and Algebraic Curves}, 2nd ed., AMS 2004, \S 15.1.1（原书此处所引）。]}

\subsection{局部 Chevalley 复形}

计算上同调 $H^\bullet(L\fg,\cA^{\fg}_{0,\mathrm{loc}})$ 的通常 Chevalley 复形是 $\cC^\bullet(L\fg,\cA^{\fg}_{0,\mathrm{loc}}) = \bigoplus_{i\ge 0}\cC^i(L\fg,\cA^{\fg}_{0,\mathrm{loc}})$，其中
\[
\cC^i(L\fg,\cA^{\fg}_{0,\mathrm{loc}}) = \Hom_{\mathrm{cont}}\Big(\bigwedge\nolimits^i(L\fg),\cA^{\fg}_{0,\mathrm{loc}}\Big),
\]
这里 $\bigwedge^i(L\fg)$ 表示拓扑向量空间 $L\fg$ 的通常第 $i$ 个外幂的自然完备化，且只考虑连续线性映射。

对 $f\in\cA^{\fg}_{0,\mathrm{loc}}$，用 $\psi^*_{b_1,-k_1}\cdots\psi^*_{b_i,-k_i}f$ 表示由下述公式定义的线性泛函
\[
\phi\in\Hom_{\mathrm{cont}}\Big(\bigwedge\nolimits^i(L\fg),\cA^{\fg}_{0,\mathrm{loc}}\Big),
\]
\begin{equation}
\phi(J^{a_1}_{m_1}\wedge\cdots\wedge J^{a_i}_{m_i}) =
\begin{cases}
(-1)^{l(\tau)}f, & ((a_1,m_1),\dots,(a_i,m_i)) = \tau((b_1,k_1),\dots,(b_i,k_i)),\\
0, & ((a_1,m_1),\dots,(a_i,m_i)) \ne \tau((b_1,k_1),\dots,(b_i,k_i)),
\end{cases}
\tag{5.6-1}
\end{equation}
其中 $\tau$ 遍历 $i$ 个字母上的对称群，$l(\tau)$ 是 $\tau$ 的长度。于是空间 $\cC^i(L\fg,\cA^{\fg}_{0,\mathrm{loc}})$ 中任一元素都可以写成这种形式的项（可能无穷）的线性组合。

微分 $d:\cC^i(L\fg,\cA^{\fg}_{0,\mathrm{loc}})\to\cC^{i+1}(L\fg,\cA^{\fg}_{0,\mathrm{loc}})$ 由公式
\[
\begin{aligned}
(d\phi)(X_1,\dots,X_{i+1}) = &\sum_{j=1}^{i+1}(-1)^{j+1}X_j\phi(X_1,\dots,\widehat{X_j},\dots,X_{i+1})\\
&+\sum_{j<k}(-1)^{j+k+1}\phi([X_j,X_k],X_1,\dots,\widehat{X_j},\dots,\widehat{X_k},\dots,X_{i+1})
\end{aligned}
\]
给出。

由顶点算子的定义可知，线性映射
\begin{equation}
\int Y(\psi^*_{a_1,n_1}\cdots\psi^*_{a_i,n_i}a^*_{\alpha_1,m_1}\cdots a^*_{\alpha_j,m_j}|0\rangle,z)\,dz,\qquad n_p\le 0,\ m_p\le 0
\tag{5.6-2}
\end{equation}
是 $\bigwedge^i(L\fg)\to\cA^{\fg}_{0,\mathrm{loc}}$ 的连续映射。这里及下文使用记号
\[
\int f(z)dz = \Res_{z=0}f(z)dz.
\]
设 $\cC^i_{\mathrm{loc}}(L\fg,\cA^{\fg}_{0,\mathrm{loc}})$ 为空间 $\Hom_{\mathrm{cont}}(\bigwedge^i(L\fg),\cA^{\fg}_{0,\mathrm{loc}})$ 中由一切形如 (5.6-2) 的线性映射张成的子空间。

\begin{lemma}
Chevalley 微分把子空间 $\cC^i_{\mathrm{loc}}(L\fg,\cA^{\fg}_{0,\mathrm{loc}})\subset\cC^i(L\fg,\cA^{\fg}_{0,\mathrm{loc}})$ 映入 $\cC^{i+1}_{\mathrm{loc}}(L\fg,\cA^{\fg}_{0,\mathrm{loc}})\subset\cC^{i+1}(L\fg,\cA^{\fg}_{0,\mathrm{loc}})$，因此
\[
\cC^\bullet_{\mathrm{loc}}(L\fg,\cA^{\fg}_{0,\mathrm{loc}}) = \bigoplus_{i\ge 0}\cC^i_{\mathrm{loc}}(L\fg,\cA^{\fg}_{0,\mathrm{loc}})
\]
是 $\cC^\bullet(L\fg,\cA^{\fg}_{0,\mathrm{loc}})$ 的子复形。
\end{lemma}

\begin{proof}
微分 $d$ 在 $\int Y(A,z)dz$ 上的作用，其中 $A$ 形如 (5.6-2)，可以写成对易子
\[
\Big[\int Q(z)dz,\int Y(A,z)dz\Big],
\]
其中 $Q(z)$ 是场
\begin{equation}
Q(z) = \sum_a J^a(z)\psi^*_a(z) - \frac{1}{2}\sum_{a,b,c}\mu^{ab}_c{:}\psi^*_a(z)\psi^*_b(z)\psi_c(z){:},
\tag{5.6-3}
\end{equation}
其中 $\mu^{ab}_c$ 表示 $\fg$ 的结构常数。因此它等于
\[
Y\Big(\int Q(z)dz\cdot A,z\Big),
\]
而这是 (5.6-2) 的形式。
\end{proof}

我们把 $\cC^i_{\mathrm{loc}}(L\fg,\cA^{\fg}_{0,\mathrm{loc}})$ 称为 $\cC^\bullet(L\fg,\cA^{\fg}_{0,\mathrm{loc}})$ 的\textbf{局部子复形}。

\begin{lemma}
$\omega$ 与 $i_*(\sigma)$ 都属于 $\cC^2_{\mathrm{loc}}(L\fg,\cA^{\fg}_{0,\mathrm{loc}})$。
\end{lemma}

\begin{proof}
先从 $\sigma$ 入手，它由公式 (5.5-5) 给出。因此上循环 $i_*(\sigma)$ 等于 $\cC^2_{\mathrm{loc}}(L\fg,\cA^{\fg}_{0,\mathrm{loc}})$ 的如下元素：
\begin{equation}
\begin{aligned}
i_*(\sigma) &= \sum_{a\le b}\sum_{n\in\mathbb{Z}}\kappa_c(J^a,J^b)n\,\psi^*_{a,-n}\psi^*_{b,n}\\
&= \sum_{a\le b}\kappa_c(J^a,J^b)\int Y(\psi^*_{a,-1}\psi^*_{b,0}|0\rangle,z)\,dz.
\end{aligned}
\tag{5.6-4}
\end{equation}

再考虑 $\omega$。结合 5.3.2 节的讨论与 5.2.5 节的公式，得到
\[
\iota(J^a_n) = \sum_{\beta\in\Delta_+}\int Y(R^\beta_a(a^*_{\alpha,0})a_{\beta,-1}|0\rangle,z)\,z^n dz,
\]
其中 $R^\beta_a$ 是多项式。于是公式 (5.5-3) 蕴含
\[
\begin{aligned}
\omega(J^a_n,J^b_m) = -\sum_{\alpha,\beta\in\Delta_+}\int\Big(&Y\Big(T\frac{\partial R^\alpha_a}{\partial a^*_{\beta,0}}\cdot\frac{\partial R^\beta_b}{\partial a^*_{\alpha,0}},w\Big)w^{n+m}\\
&+nY\Big(\frac{\partial R^\alpha_a}{\partial a^*_{\beta,0}}\frac{\partial R^\beta_b}{\partial a^*_{\alpha,0}},w\Big)w^{n+m-1}\Big)\,dw.
\end{aligned}
\]
因此
\begin{equation}
\begin{aligned}
\omega = -\sum_{a\le b;\,\alpha,\beta\in\Delta_+}\int\Big(&Y\Big(\psi^*_{a,0}\psi^*_{b,0}T\frac{\partial R^\alpha_a}{\partial a^*_{\beta,0}}\cdot\frac{\partial R^\beta_b}{\partial a^*_{\alpha,0}},z\Big)\\
&+Y\Big(\psi^*_{a,-1}\psi^*_{b,0}\frac{\partial R^\alpha_a}{\partial a^*_{\beta,0}}\frac{\partial R^\beta_b}{\partial a^*_{\alpha,0}},z\Big)\Big)\,dz
\end{aligned}
\tag{5.6-5}
\end{equation}
（最后两个公式中省略了 $|0\rangle$）。所以它属于空间 $\cC^2_{\mathrm{loc}}(L\fg,\cA^{\fg}_{0,\mathrm{loc}})$。
\end{proof}

% [Tutor Note: (5.6-3) 的 Q(z) 值得留意：它正是 Chevalley 微分在顶点代数语言中的“内蕴形式”，形如 BRST 流。引理 5.6.1 的实质是 d(∫Y(A,z)dz) = [∫Q dz, ∫Y(A,z)dz]，即微分等于与 Q 的对易子；这条恒等式后面在引理 5.6.3 中再次使用，是整节的机械核心。]

需要证明上循环 $i_*(\sigma)$ 与 $\omega$ 在局部复形 $\cC^\bullet_{\mathrm{loc}}(L\fg,\cA^{\fg}_{0,\mathrm{loc}})$ 中代表同一上同调类。我们将证明，这等价于检验这两个上循环在李子代数 $L\fh\subset L\fg$ 上的限制相同。这些限制容易计算，而且它们确实相等。过渡到 $L\fh$ 靠的是 Shapiro 引理的一个版本，下一小节将予以解释。

\subsection{另一个复形}

给定李代数 $\mathfrak{l}$，把 $L\mathfrak{l} = \mathfrak{l}((t))$ 的李子代数 $\mathfrak{l}[[t]]$ 记作 $L_+\mathfrak{l}$。

注意，李代数 $L_+\fg$ 自然地作用在空间
\[
M_{\fg,+} = \C[a^*_{\alpha,n}]_{\alpha\in\Delta_+,\,n\le 0}\simeq \C[y_{\alpha,n}]_{\alpha\in\Delta_+,\,n\ge 0}
\]
上，该空间与 $\cU[[t]]\simeq N_+[[t]]$ 上函数环等同。把标准 Chevalley 复形
\[
\cC^\bullet(L_+\fg,M_{\fg,+}) = \Hom_{\mathrm{cont}}\Big(\bigwedge\nolimits^\bullet L_+\fg,M_{\fg,+}\Big)
\]
与张量积 $M_{\fg,+}\otimes\bigwedge^\bullet_{\fg,+}$ 等同，其中
\[
\bigwedge\nolimits^\bullet_{\fg,+} = \bigwedge(\psi^*_{a,n})_{n\le 0}.
\]
在 $\cC^\bullet(L_+\fg,M_{\fg,+})$ 上引入超导子 $T$，其作用为
\[
T\cdot a^*_{a,n} = -(n-1)a^*_{a,n-1},\qquad
T\cdot\psi^*_{a,n} = -(n-1)\psi^*_{a,n-1}.
\]
% [Tutor Note: 此处印作 a^*_{a,n}，与上一行 M_{g,+} = C[a^*_{α,n}] 的指标 α 不一致，疑为 a^*_{α,n} 之误；译文保留印出的 a 以便对照，登记于报告。]

有线性映射
\[
\int:\cC^\bullet(L_+\fg,M_{\fg,+})\to\cC^\bullet_{\mathrm{loc}}(L\fg,\cA^{\fg}_{0,\mathrm{loc}}),
\]
它把 $A\in\cC^\bullet(L_+\fg,M_{\fg,+})$ 送到 $\int Y(A,z)dz$（回忆 $\int$ 取出形式幂级数的第 $(-1)$ 个 Fourier 系数）。\transnote{卡点：上文 $T\cdot a^*_{a,n}=-(n-1)a^*_{a,n-1}$ 的下标印作罗马字母 $a$，而本节 $M_{\fg,+}=\C[a^*_{\alpha,n}]_{\alpha\in\Delta_+,n\le0}$ 用的是根指标 $\alpha$——同一字母在两处指代不同对象。背景：此处的 $T$ 是复形 $\cC^\bullet(L_+\fg,M_{\fg,+})$ 上的超导子，它应把 $a^*_{\alpha,n}$ 送到 $-(n-1)a^*_{\alpha,n-1}$（与 5.4.1 节 $\deg a^*_{\alpha,n}=-n$ 的约定相容），对 $\psi^*_{a,n}$ 的作用同理，那里的 $a$ 才是 $\fg$ 的基指标。结论：原书 p.156 的 $a$ 系 $\alpha$ 之误，译文照印；读者按 $a^*_{\alpha,n}$ 理解即可。[来源：原书 p.156 同页 $M_{\fg,+}$ 的定义与 p.142 的 $[T,a^*_{\alpha,n}]=-(n-1)a^*_{\alpha,n-1}$。]}

回忆在任何顶点代数 $V$ 中都有恒等式 $Y(TA,z) = \partial_zY(A,z)$。因此若 $A\in\Im T$，则 $\int Y(A,z)dz = 0$。

\begin{lemma}
映射 $\int$ 定义同构
\[
\cC^\bullet_{\mathrm{loc}}(L\fg,\cA^{\fg}_{0,\mathrm{loc}})\simeq \cC^\bullet(L_+\fg,M_{\fg,+})/(\Im T+\C),
\]
且 $\operatorname{Ker}T = \C$。此外，下图交换：
\[
\begin{CD}
\cC^\bullet(L_+\fg,M_{\fg,+}) @>T>> \cC^\bullet(L_+\fg,M_{\fg,+}) @>\int>> \cC^\bullet_{\mathrm{loc}}(L\fg,\cA^{\fg}_{0,\mathrm{loc}})\\
@AdAA @AdAA @AdAA\\
\cC^\bullet(L_+\fg,M_{\fg,+}) @>T>> \cC^\bullet(L_+\fg,M_{\fg,+}) @>\int>> \cC^\bullet_{\mathrm{loc}}(L\fg,\cA^{\fg}_{0,\mathrm{loc}}).
\end{CD}
\]
\end{lemma}

\begin{proof}
容易看出，Chevalley 复形 $\cC^\bullet(L_+\fg,M_{\fg,+})$ 的微分按公式 $A\mapsto\int Q(z)dz\cdot A$ 作用，其中 $Q(z)$ 由公式 (5.6-3) 给出。于是本引理可用引理 5.6.1 证明中的论证推出。
\end{proof}

考虑双复形
\[
\begin{CD}
\C @>>> \cC^\bullet(L_+\fg,M_{\fg,+}) @>T>> \cC^\bullet(L_+\fg,M_{\fg,+}) @>>> \C\\
@. @AdAA @AdAA @.\\
\C @>>> \cC^\bullet(L_+\fg,M_{\fg,+}) @>T>> \cC^\bullet(L_+\fg,M_{\fg,+}) @>>> \C.
\end{CD}
\]
按引理，复形 $\cC^\bullet_{\mathrm{loc}}(L\fg,\cA^{\fg}_{0,\mathrm{loc}})$ 的上同调由谱序列的第二项给出，其中第零个微分是竖直的。

先从计算该谱序列的第一项入手。注意李代数 $L_+\fn_+$ 带有 $Q\times\mathbb{Z}_+$ 分次，其中 $Q$ 是 $\fg$ 的根格，分次由公式 $\deg e_{\alpha,n} = (\alpha,n)$ 定义。因此泛包络代数 $U(L_+\fn_+)$ 也是 $\mathbb{Z}_+\times Q$-分次的。此外，对 $\gamma\in Q$、$n\in\mathbb{Z}_+$，$U(L_+\fn_+)$ 的所有齐次分量 $U(L_+\fn_+){}_{(\gamma,n)}$ 都是有限维的。相应的限制对偶空间是
\[
U(L_+\fn_+)^\vee = \bigoplus_{(\gamma,n)\in Q\times\mathbb{Z}_+}\big(U(L_+\fn_+)_{(\gamma,n)}\big)^*.
\]
它带有自然的 $L_+\fn_+$-模结构。

现在令
\[
\operatorname{Coind}_{L_+\fb_-}^{L_+\fg}\C \overset{\mathrm{def}}{=}\Hom^{\mathrm{res}}_{U(L_+\fb_-)}(U(L_+\fg),\C)
\]
（比较公式 (5.2-9)）。右端由一切“受限”线性泛函 $U(L_+\fg)\to\C$ 组成，它们在 $U(\fb_-)$ 对 $U(\fg)$ 的左作用之下不变；这里“受限”是指只考虑这样的泛函：它们是支撑在 $U(\fg)$ 的直和项 $U(\fb_-)\otimes U(\fn_+){}_{(\gamma,n)}$ 上的映射的有限线性组合（相对于向量空间同构 $U(\fg)\simeq U(\fb_-)\otimes U(\fn_+)$）。换言之，作为 $L_+\fn_+$-模，$\operatorname{Coind}_{L_+\fb_-}^{L_+\fg}\C\simeq U(L_+\fn_+)^\vee$。\transnote{卡点：这一段里的 $U(\fb_-)$、$U(\fn_+)$、$U(\fg)$ 是 $U(L_+\fb_-)$、$U(L_+\fn_+)$、$U(L_+\fg)$ 的简写（原书在此省略下标 $+$），不能读成有限维对象；另外这里为何一步就能把复形换小。背景：先由引理（$M_{\fg,+}\simeq\operatorname{Coind}_{L_+\fb_-}^{L_+\fg}\C$）用 Shapiro 引理把 $\cC^\bullet(L_+\fg,M_{\fg,+})$ 换成 $\cC^\bullet(L_+\fb_-,\C)$，再用 $L_+\fb_-=L_+\fn_-\oplus L_+\fh$ 的 Serre--Hochschild 谱序列；$L_+\fh$ 的对角作用给出内分次，最终只剩 $H^\bullet(L_+\fn_-,\C)$ 需要计算。结论：这条“Shapiro $\to$ Serre--Hochschild $\to$ 只剩幂幺部分”的路线是 5.6.4 节的全部技术内容。[来源：原书 p.158--159；Fuks, \emph{Cohomology of Infinite-Dimensional Lie Algebras}, Consultants Bureau 1986, \S 5.1--5.4（原书此处所引）。]}

% [Tutor Note: 这一段里 U(b_−)、U(n_+)、U(g) 是 U(L_+b_−)、U(L_+n_+)、U(L_+g) 的简写（原书在此省略了下标 +），不要误读成有限维对象。余诱导模与限制对偶 U(L_+n_+)^∨ 的等同是下面 Shapiro 引理的入口。]

\begin{lemma}
$L_+\fg$-模 $M_{\fg,+}$ 同构于余诱导模 $\operatorname{Coind}_{L_+\fb_-}^{L_+\fg}\C$。
\end{lemma}

\begin{proof}
逐字照搬我们在 5.2.3 与 5.2.4 节给出的、关于 $\operatorname{Fun}N_+$ 作为 $\fg$-模同构于 $\operatorname{Coind}_{\fb}^{\fg}\C$ 的证明。先证明 $M_{\fg,+} = \operatorname{Fun}L_+N_+$ 作为 $L_+\fn_+$-模同构于 $U(L_+\fn_+)^\vee$：为此证明 $M_{\fg,+}$ 的单项式基与 $U(L_+\fn_+)$ 的 PBW 单项式基关于两个模的一个 $L_+\fn_+$-不变配对互为对偶。然后利用这一事实证明自然同态 $M_{\fg,+}\to\operatorname{Coind}_{L_+\fb_-}^{L_+\fg}\C$ 是同构。
\end{proof}

定义复形映射
\[
\mu':\cC^\bullet(L_+\fg,M_{\fg,+})\to\cC^\bullet(L_+\fb_-,\C)
\]
如下：若 $\gamma$ 是复形
\[
\cC^\bullet(L_+\fg,M_{\fg,+}) = \Hom_{\mathrm{cont}}\Big(\bigwedge\nolimits^i L_+\fg,M_{\fg,+}\Big)
\]
中的一个 $i$-上链，则 $\mu'(\gamma)$ 按定义是 $\gamma$ 在 $\bigwedge^i L_+\fb_-$ 上的限制与自然投影 $M_{\fg,+} = \operatorname{Coind}_{L_+\fb_-}^{L_+\fg}\C\to\C$ 的复合。显然 $\mu'$ 是复形态射。下面是 Shapiro 引理的一个例子（证明见 [Fuks 1986] 第 5.4 节）。

\begin{lemma}
映射 $\mu'$ 在上同调层面诱导同构，即
\begin{equation}
H^\bullet(L_+\fg,M_{\fg,+})\simeq H^\bullet(L_+\fg,\operatorname{Coind}_{L_+\fb_-}^{L_+\fg}\C)\simeq H^\bullet(L_+\fb_-,\C).
\tag{5.6-6}
\end{equation}
\end{lemma}

现在计算 (5.6-6) 的右端。由于 $L_+\fn_-\subset L_+\fb_-$ 是理想，且 $L_+\fb_-/L_+\fn_-\simeq L_+\fh$，由 Serre--Hochschild 谱序列（见 [Fuks 1986] 第 5.1 节）得到
\[
H^n(L_+\fb_-,\C) = \bigoplus_{p+q=n}H^p(L_+\fh,H^q(L_+\fn_-,\C)).
\]
但 $\fh\otimes 1\in L_+\fh$ 对角地作用在 $H^\bullet(L_+\fn_-,\C)$ 上，诱导出一个内分次。按 [Fuks 1986] 第 5.2 节，
\[
H^p(L_+\fh,H^q(L_+\fn_-,\C)) = H^p(L_+\fh,H^q(L_+\fn_-,\C)_0),
\]
其中 $H^q(L_+\fn_-,\C)_0$ 是 $\fh\otimes 1$ 以 0 作用的子空间。显然，$H^0(L_+\fn_-,\C)_0$ 是标量构成的一维子空间 $\C\subset H^0(L_+\fn_-,\C)$，而当 $q\ne 0$ 时 $H^q(L_+\fn_-,\C)_0 = 0$。于是得到
\[
H^p(L_+\fh,H^q(L_+\fn_-,\C)) = H^p(L_+\fh,\C).
\]
此外还有如下结果。定义复形映射
\begin{equation}
\mu:\cC^\bullet(L_+\fg,M_{\fg,+})\to\cC^\bullet(L_+\fh,\C)
\tag{5.6-7}
\end{equation}
如下：若 $\gamma$ 是复形
\[
\cC^\bullet(L_+\fg,M_{\fg,+}) = \Hom_{\mathrm{cont}}\Big(\bigwedge\nolimits^i L_+\fg,M_{\fg,+}\Big)
\]
中的一个 $i$-上链，则 $\mu(\gamma)$ 按定义是 $\gamma$ 在 $\bigwedge^i L_+\fh$ 上的限制与自然投影
\begin{equation}
p:M_{\fg,+} = \operatorname{Coind}_{L_+\fb_-}^{L_+\fg}\C\to\C
\tag{5.6-8}
\end{equation}
的复合。

\begin{lemma}
映射 $\mu$ 在上同调层面诱导同构，即
\[
H^\bullet(L_+\fg,M_{\fg,+})\simeq H^\bullet(L_+\fh,\C).
\]
特别地，复形 $\cC^i(L_+\fg,M_{\fg,+})$ 中一个上循环的上同调类，由它在 $\bigwedge^i L_+\fh$ 上的限制唯一决定。
\end{lemma}

\begin{proof}
由 Serre--Hochschild 谱序列的构造（见 [Fuks 1986] 第 5.1 节）可知，限制映射 $\cC^\bullet(L_+\fb_-,\C)\to\cC^\bullet(L_+\fh,\C)$ 在上同调层面诱导同构。把这一点与引理 5.6.5 结合，便得到本引理的结论。
\end{proof}

由于 $L_+\fh$ 是阿贝尔的，有
\[
H^\bullet(L_+\fh,\C) = \bigwedge\nolimits^\bullet(L_+\fh)^*,
\]
于是
\[
H^\bullet(L_+\fg,M_{\fg,+})\simeq \bigwedge\nolimits^\bullet(L_+\fh)^*.
\]
显然这个同构与 $T$ 在两侧的作用相容。$T$ 在右端作用的核等于标量子空间。因此得到
\[
H^\bullet_{\mathrm{loc}}(L\fg,\cA^{\fg}_{0,\mathrm{loc}})\simeq \bigwedge\nolimits^\bullet(L_+\fh)/(\Im T+\C).
\]

% [Tutor Note: 这一串化简（Shapiro 引理 → Serre–Hochschild 谱序列 → 只剩 Cartan）是全节技术核心：它说明局部复形的上同调类完全由其在 ⋀^•(L_+h) 上的限制决定。有了这句话，定理 5.6.8 就只需在 Cartan 上做一次初等计算——这正是 5.6.4 的引理 5.6.7 所表述的内容。]

\subsection{限制上循环}

$\cC^i_{\mathrm{loc}}(L\fg,\cA^{\fg}_{0,\mathrm{loc}})$ 中任一上循环都是复形
\[
\cC^i(L\fg,\cA^{\fg}_{0,\mathrm{loc}}) = \Hom_{\mathrm{cont}}\Big(\bigwedge\nolimits^i(L\fg),\cA^{\fg}_{0,\mathrm{loc}}\Big)
\]
中的上循环，因此可以限制到 $\bigwedge^i(L\fh)$ 上。

下面这个引理是关键结果，它使我们能证明 $\omega$ 与 $i_*(\sigma)$ 定义同一上同调类。

\begin{lemma}
$\cC^i_{\mathrm{loc}}(L\fg,\cA^{\fg}_{0,\mathrm{loc}})$ 中任意两个在 $\bigwedge^i(L\fh)$ 上限制相同的上循环，代表同一上同调类。
\end{lemma}

\begin{proof}
需要证明：$\cC^i_{\mathrm{loc}}(L\fg,\cA^{\fg}_{0,\mathrm{loc}})$ 中任一在 $\bigwedge^i(L\fh)$ 上限制为零的上循环 $\phi$，都等价于零上循环。按上面的计算，$\phi$ 可以写成 $\int Y(A,z)dz$，其中 $A$ 是 $C^i(L_+\fg,M_{\fg,+}) = \Hom_{\mathrm{cont}}(\bigwedge^i L_+\fg,M_{\fg,+})$ 中的上循环。但此时 $A$ 到 $\bigwedge^i L_+\fh\subset\bigwedge^i L_+\fg$ 的限制（记作 $\bar A$）必属于算子 $T$ 的像，于是 $p(A) = \mu(A)\in\bigwedge^i(L_+\fh)^*$ 也属于 $T$ 的像（这里 $p$ 是 (5.6-8) 中定义的投影，$\mu$ 是 (5.6-7) 中定义的映射）。于是 $\mu(A) = T(h)$，其中 $h\in\bigwedge^i(L_+\fh)^*$。由于 $T$ 与微分对易，且 $T$ 在 $\bigwedge^i(L_+\fh)^*$ 上的核由标量构成，可知 $h$ 也是 $C^i(L_+\fh,\C) = \bigwedge^i(L_+\fh)^*$ 中的上循环。

按引理 5.6.6，映射 $\mu$ 在上同调上诱导同构。因此 $h = \mu(B)$，其中 $B$ 是 $C^i(L_+\fg,M_{\fg,+})$ 中的某个上循环。由定义显然可见 $\mu$ 与平移算子 $T$ 的作用对易。所以上循环 $A$ 与 $T(B)$ 在 $C^i(L_+\fg,M_{\fg,+})$ 中等价。于是 $\phi$ 等价于零上循环。
\end{proof}

现在可以证明本章的主要结果了。

\begin{theorem}
上循环 $\omega$ 与 $i_*(\sigma)$ 在 $H^2_{\mathrm{loc}}(L\fg,\cA^{\fg}_{0,\mathrm{loc}})$ 中代表同一上同调类。因此存在同态 $L\fg\to\cT^{\fg}_{\mathrm{loc}}$ 的提升，得到同态 $\gghat_{\kappa_c}\to\cA^{\fg}_{\le1,\mathrm{loc}}$，使得 $\mathbf{1}\mapsto 1$。此外，这个同态可以选得使
\begin{equation}
J^a(z)\mapsto Y(P_a(a^*_{\alpha,0},a_{\beta,-1})|0\rangle,z) + Y(B_a,z),
\tag{5.6-9}
\end{equation}
其中 $P_a$ 是公式 (5.2-11) 中引入的多项式，$B_a$ 是 $a^*_{\alpha,n}$（$n\le 0$）的多项式，次数为 1（相对于分次 $\deg a^*_{\alpha,n} = -n$）。
\end{theorem}

\begin{proof}
按引理 5.6.7，只需检验 $\omega$ 与 $i_*(\sigma)$ 到 $\bigwedge^2(L\fh)$ 上的限制相同。对一切 $h,h'\in\fh$ 有
\[
i_*(\sigma)(h_n,h'_m) = n\kappa_c(h,h')\delta_{n,-m}.
\]
现在计算 $\omega$ 的限制。由 5.2.5 节与 5.3.2 节给出的公式得到
\begin{equation}
\iota(h(z)) = -\sum_{\alpha\in\Delta_+}\alpha(h)\,{:}a^*_\alpha(z)a_\alpha(z){:}
\tag{5.6-10}
\end{equation}
（回忆 $h(z) = \sum_{n\in\mathbb{Z}}h_nz^{-n-1}$）。因此，与 5.5.3 节末尾的计算相同，得到
\[
\omega(h_n,h'_m) = -n\delta_{n,-m}\sum_{\alpha\in\Delta_+}\alpha(h)\alpha(h').
\]
关键的观察是
\begin{equation}
\sum_{\alpha\in\Delta_+}\alpha(h)\alpha(h') = \kappa_c(h,h'),
\tag{5.6-11}
\end{equation}
因为按定义 $\kappa_c(\cdot,\cdot) = -\tfrac12\kappa_{\fg}(\cdot,\cdot)$，而对 Killing 型 $\kappa_{\fg}$ 有
\[
\kappa_{\fg}(h,h') = 2\sum_{\alpha\in\Delta_+}\alpha(h)\alpha(h').
\]
% [Tutor Note: (5.6-11) 的符号与紧接着给出的理由矛盾：由 κ_c = −(1/2)κ_g 与 κ_g(h,h′) = 2Σα(h)α(h′)，应得 Σα(h)α(h′) = −κ_c(h,h′)，即 (5.6-11) 右端应带负号。只有取 −κ_c，下文 ω(h_n,h′_m) = nκ_c(h,h′)δ_{n,−m} 才与 i_*(σ) 一致；译文保留印出的形式，已登记为疑似笔误。]
因此得到
\[
\omega(h_n,h'_m) = n\kappa_c(h,h')\delta_{n,-m},
\]
于是上循环 $\omega$ 与 $i_*(\sigma)$ 到 $\bigwedge^2(L\fh)$ 的限制相同。按引理 5.6.7，它们代表 $H^2_{\mathrm{loc}}(L\fg,\cA^{\fg}_{0,\mathrm{loc}})$ 中的同一上同调类。\transnote{卡点：(5.6-11) 的右端少了负号。背景：由 $\kappa_c=-\tfrac12\kappa_{\fg}$ 与紧接其后的 $\kappa_{\fg}(h,h')=2\sum_{\alpha}\alpha(h)\alpha(h')$ 立得 $\sum_{\alpha\in\Delta_+}\alpha(h)\alpha(h')=\tfrac12\kappa_{\fg}(h,h')=-\kappa_c(h,h')$，而原书印作 $=\kappa_c(h,h')$。结论：只有把它读作 $-\kappa_c$，下一行 $\omega(h_n,h'_m)=n\kappa_c(h,h')\delta_{n,-m}$ 才与 $i_*(\sigma)$ 的同一式相符——而这一步正是证明的关键。译文照原书保留印出的形式。[来源：原书 p.161 同页（$\kappa_c$、$\kappa_{\fg}$ 两式）与紧随的结论式（PDF p.174）。]}

因此存在 $\gamma\in C^1_{\mathrm{loc}}(L\fg,\cA^{\fg}_{0,\mathrm{loc}})$，使得 $\omega+d\gamma = i_*(\sigma)$。可以把 $\gamma$ 写成
\[
\gamma = \sum_a\int\psi^*_a(z)Y(B_a,z)\,dz,\qquad B_a\in M_{\fg,+}.
\]
由引理 5.6.2 证明中的计算可知，$\omega$ 与 $i_*(\sigma)$ 都是 $C^2_{\mathrm{loc}}(L\fg,\cA^{\fg}_{0,\mathrm{loc}})$ 中次数为 0 的齐次元素。因此 $\gamma$ 可以选为次数 0，即 $B_a$ 可选为次数 1。

于是由引理 5.5.3，公式 (5.6-9) 定义了形如 (5.6-9) 的李代数同态 $\gghat_{\kappa_c}\to\cA^{\fg}_{\le1,\mathrm{loc}}$，且 $\mathbf{1}\mapsto 1$。证明完毕。
\end{proof}

% [Tutor Note: 注意定理 5.6.8 只断言“存在”这样的提升，而 (5.6-9) 给出了它的标准形：主项 Y(P_a(a^*_{α,0},a_{β,−1})|0⟩,z) 是正常的自由场实现，修正项 B_a 是 a^*_{α,n} 的一次多项式。第 6 章把这一同态表述为顶点代数同态，即 Wakimoto 模；B_a 的显式计算是 6.2 节“形变到其他水平”的起点。]

\subsection{其他旗簇的障碍}

上面证明中用到的一个关键事实是：公式 (5.6-11) 左端定义的 $\fh$ 上内积，是 $\fg$ 上某个不变内积在 $\fh$ 上的限制。我们随后发现，这个不变内积正是临界内积 $\kappa_c$。

注意，$\fh$ 上一个内积等于 $\fg$ 上某不变内积在 $\fh$ 上的限制，其充要条件是它 $W$-不变，其中 $W$ 是 $\fg$ 的外尔群，作用在 $\fh$ 上。对公式 (5.6-11) 左端定义的内积，这一点显然成立，因为集合 $\Delta_+$ 是 $W$-不变的。但这一观察同时表明：在 $G$ 的所有旗簇中，只有“全”旗簇 $G/B_-$ 能使上述上同调障碍被克服。

确实，考虑更一般的旗簇 $G/P$，其中 $P$ 是抛物子群。例如 $\fg=\fsl_n$ 时，抛物子群的共轭类对应 $n$ 的非平凡分拆。对应分拆 $n = n_1+\cdots+n_k$ 的标准（下）抛物子群由块大小为 $n_1,\dots,n_k$ 的“块下三角”矩阵构成。相应的旗簇是旗 $V_1\subset\cdots\subset V_{k-1}\subset\C^n$ 构成的簇，其中 $\dim V_i = n_1+\cdots+n_i$。

设 $P$ 是异于 Borel 子群的抛物子群，$\fp$ 是 $\fg$ 中相应的李子代数。不妨设 $\fp$ 包含 $\fb_-$，且在 $H$ 的伴随作用下不变。于是有分解 $\fg = \fp\oplus\mathfrak{n}$，其中 $\mathfrak{n}$ 是 $\fn_+$ 的李子代数，由生成元 $e_\alpha$ 张成，$\alpha$ 遍历真子集 $\Delta_P\subset\Delta_+$。对 $G/P$ 可以展开与 $G/B_-$ 相同的形式理论：李子代数 $\mathfrak{n}$ 现在扮演 $\fn_+$ 的角色，其李群同构于 $G/P$ 的一个开稠子集，等等。特别地，公式 (5.6-10) 被修改为：求和只遍历 $\Delta_S$，而不是遍历 $\Delta_+$。于是在相应的上循环中，得到由公式
\[
\nu_P(h,h') = \sum_{\alpha\in\Delta_P}\alpha(h)\alpha(h')
\]
给出的 $\fh$ 上内积。但 $\Delta_+$ 的任何真子集都不是 $W$-不变的。因此内积 $\nu_P$ 不是 $W$-不变的，从而不能作为 $\fg$ 上不变内积在 $\fh$ 上的限制。于是我们的论证失效；这说明当抛物子群 $P$ 不是 Borel 子群时，无法把从 $L\fg$ 到相应向量场李代数的同态，提升为从中心扩张 $\gghat_\kappa$ 到微分算子代数的同态。\transnote{卡点：同段先说 $\alpha$ 遍历真子集 $\Delta_P\subset\Delta_+$，紧接着又说“求和只遍历 $\Delta_S$”，末了又写 $\nu_P(h,h')=\sum_{\alpha\in\Delta_P}$——$\Delta_S$ 与 $\Delta_P$ 混用。背景：二者指同一个集合，即由 Levi 因子 $\mathfrak{m}$ 对应的单根子集 $S$ 生成的那些正根；原书 p.162 同页三处并用两种记号。结论：应统一读作 $\Delta_P$；译文照原书保留印出的 $\Delta_S$。正是这一“缩小后的根集”使 $\nu_P$ 失去 $W$-不变性，从而不能成为 $\fg$ 上不变内积的限制。[来源：原书 p.162 同页三处（PDF p.175）。]}
% [Tutor Note: 这里 Δ_S 与上一句的 Δ_P 指同一个集合（由 Levi 因子对应的单根子集 S 生成的正根），原书两种记号混用；可统一理解为 Δ_P。另注意 ν_P 与 (5.6-11) 的差别仅在于求和范围，正是这个“缺少的根”使 ν_P 失去 W-不变性，从而无法成为 g 上不变内积的限制。]

这一点体现了全旗簇 $G/B_-$ 的特殊地位。它与如下事实密切相关：$G/B_-$ 的第一 Pontryagin 类为零，而对其他旗簇它非零，从而产生“反常”。不过，我们将在 6.3 节说明一般情形下如何克服这一反常。基本想法是添加 Levi 子群 $\mathfrak{m}$（$\fp$ 的 Levi 因子）的仿射化的额外表示。适当选取这些表示的水平，就能抵消反常。

% [Tutor Note: “第一 Pontryagin 类为零”是几何侧对同一上同调障碍的解释：G/B 的切丛的第一 Pontryagin 类正比于临界水平的上同调类，故只有全旗簇没有反常。6.3 节的半无限抛物诱导正是用 Levi 子群的仿射化表示的额外因子来补偿这个类，这是第 6 章与第 9 章的技术预演。]

总结本章结果：我们现已得到从 $\gghat_{\kappa_c}$ 到李代数 $\cA^{\fg}_{\le1,\mathrm{loc}}$ 的“自由场实现”同态，后者由对应于全旗簇 $G/B_-$ 上微分算子的“自由场” $a_\alpha(z)$ 与 $a^*_\alpha(z)$ 所构造的顶点算子的 Fourier 系数组成。下一章我们将在顶点代数语言中解释这个同态。
